\begin{abstract}
Questo lavoro presenta \textit{GraphMate}, un sistema per la predizione di mosse
tattiche negli scacchi basato su \textit{Graph Neural Networks} (GNN). L'obiettivo
del progetto è studiare l'efficacia di una rappresentazione strutturata della
scacchiera, nella quale una posizione non viene trattata come una semplice matrice
di celle, ma come un grafo capace di descrivere esplicitamente le relazioni
scacchistiche tra le case e i pezzi.

Il dataset principale è costruito a partire dai puzzle pubblici di Lichess,
selezionando problemi di matto forzato da \textit{mate-in-1} a
\textit{mate-in-5}. Dopo le operazioni di filtraggio, validazione,
deduplicazione e bilanciamento, ciascun problema viene formulato come un task
supervisionato di predizione della prima mossa risolutiva. Ogni posizione viene
quindi trasformata in un grafo diretto composto da 64 nodi, corrispondenti alle
case della scacchiera, e da archi che rappresentano relazioni quali mosse legali,
attacchi, difese, inchiodature e scacchi. Ai nodi e agli archi sono associate
feature locali, integrate da feature globali che descrivono lo stato complessivo
della posizione.

Su questa rappresentazione viene sviluppata e confrontata una famiglia di modelli
basati su un encoder \textit{Graph Attention Network}. La baseline formula il
problema come classificazione su un vocabolario globale di mosse UCI; successive
varianti introducono il mascheramento delle mosse illegali e, soprattutto, uno
scoring diretto delle sole mosse legali disponibili nella posizione. Vengono
inoltre analizzate due estensioni: l'integrazione di informazioni temporali
sintetiche e un'architettura di \textit{reranking} a due stadi, nella quale le
migliori candidate proposte dal modello vengono rivalutate considerando
esplicitamente la posizione risultante dopo la loro applicazione.

La valutazione comprende sia dati derivati dalla distribuzione Lichess sia un
benchmark esterno di composizioni scacchistiche classiche, utilizzato per
analizzare la capacità di generalizzazione dei modelli. Su quest'ultimo benchmark,
il modello con reranking post-mossa raggiunge un'accuratezza Top-1 del 15,5\%,
rispetto all'11,0\% del modello candidate-based di base, mostrando il beneficio
dell'analisi esplicita dello stato risultante dalla mossa. Come ulteriore termine
di confronto vengono valutati modelli linguistici locali mediante un protocollo
rigoroso di generazione di una singola mossa UCI; in tale configurazione le
prestazioni risultano nettamente inferiori a quelle dei modelli specializzati.

I risultati evidenziano come l'integrazione della struttura del dominio nello
spazio decisionale, attraverso la rappresentazione a grafo, il vincolo esplicito
di legalità e l'analisi delle posizioni successive alla mossa, costituisca un
approccio promettente per la modellazione neurale di problemi tattici negli
scacchi. Il progetto fornisce inoltre una pipeline riproducibile che comprende
preprocessing dei dati, costruzione dei grafi, addestramento, gestione dei
checkpoint e protocolli di valutazione.
\end{abstract}